% page 370 of the facsimile (image p-369 of the scan)
% block 154.9 x 233.3 mm ; left margin 8.0 mm
\pgc{-12.700mm}{0.000mm}{
\l{\cel{3}362}
\l{}
\l{\sou{meandro}. (\sou{arkitekt. Greka}). Ornamento-desegnuro, formacita}
\l{\cel{3}ek linei e vergeti qui interkrucumas. - F.}
\l{}
\l{\sou{meceno}. Personego richa e socio-povoza, qua favoras peku-}
\l{\cel{3}niale la literaturisti e la artisti. - DEFIRS.}
\l{}
\l{\sou{mecho}. Kordono, bendo, asembluro de kotono e kanabo, cir-}
\l{\cel{3}kumata da sebo o de vaxo, por facar kandeli, bujii, ceri,}
\l{\cel{3}o -imbibite ye oleo, petrolo, e c. - por brular en lampo.}
\l{}
\l{}
\l{\sou{medalio}. I. Peco de metalo, estampita por perpetuigar la}
\l{\cel{3}memoro di ago remarkinda, di personego famoza o glorioza,}
\l{\cel{3}e qua reprezentas objekto o legendo qua koncernas lu. -}
\l{\cel{3}II. Peco de metalo qua reprezentas Iesu, la Virgino, e c.,}
\l{\cel{3}od irgaltra objekto di santeso,quan la katoliki +weras}
\l{\cel{3}sur su od igas pendar de sua chapleto. - III. Peco de me-}
\l{\cel{3}talo, donita premie a la laureato di ula konkursi - rekom-}
\l{\cel{3}penso a la personi qui agis prodi su-sakrifikala - memor-}
\l{\cel{3}igive, a ti qui partoprenis ula kampanio od ecelis per}
\l{\cel{3}milit-agi. IV. Peco de metalo, kun numero, quan (en ula}
\l{\cel{3}landi) devas +werar la personi qui esas portisti di kar-}
\l{\cel{3}gaji, e c. - DEFIRS.}
\l{}
\l{\sou{medaliono}. Juvelo plata, ronda od ovala, en qua onu enklozas}
\l{\cel{3}portreto, haro-lokli, reliquii, e c. - DEFIRS.}
\l{}
\l{\sou{medio}. I. To quo cirkondas omna-sinse. - II. Ensemblo ek la}
\l{\cel{3}socio, mori, eventaji, en qui vivas o vivis persono, od}
\l{\cel{3}en qui eventas od eventis fakto. -}
\l{}
\l{\sou{mediacar}. (\sou{netrans., pri)}. Agar por konciliar personi qui}
\l{\cel{3}deskonkordas (anke pri aferi komercala). - DEFIS.}
\l{}
\l{\sou{mediano.}\cel{2}(\sou{geom.)}\cel{2}(\sou{pri figuro plana)}. Lineo rekta qua divi-}
\l{\cel{3}das aden du parti interegala, omna kordi, paralela ye un}
\l{\cel{3}direciono. - Rekto qua juntas la somito di triangulo a}
\l{\cel{3}le mezo di la latero opozita. - DEFI.}
\l{}
\l{\sou{mediat}a. I. (\sou{aludante la kozi}) Qua, situita inter du termini,}
\l{\cel{3}efikas por pasigar de un de la termini a la altra. - II.}
\l{\cel{3}(\sou{aludante la personi}). Dicesas pri la ago da persono quan}
\l{\cel{3}onu uzas por atingar ta o ca rezultajo.DEFIS.}
\l{}
\l{\sou{medicino}. La arto risanigar, fondita sur la cienco pri la}
\l{\cel{3}morbi e la medikamenti. - DEFIRS.}
\l{}
\l{\sou{mediko.}\cel{2}Persono a qua esas grantita la diplomo \textquotedbl{}doktoro pri}
\l{\cel{3}medicino\textquotedbl{}. - DEFIS.}
\l{}
\l{\sou{medikamento}. Substanco uzata, interne od extere, kom remedi-}
\l{\cel{3}ilo lor maladeso. - DEFIS.}
}
